\section{Структура на проекта, конфигурация и експлоатационна среда}

Освен архитектурата и домейн логиката, всяко завършено приложение следва да бъде разглеждано и от operational перспектива: как е организирано в хранилището, какви зависимости използва, как се компилира, къде се съхраняват ресурсите и какви артефакти се генерират в резултат от build процеса. За настоящата система това е особено важно, тъй като проектът използва класическа Maven организация, комбинирана със Spring Boot конфигурация, MySQL свързаност, SMTP настройка и набор от HTML ресурси.

Тази глава анализира именно този аспект на разработката и показва как логическата структура на приложението се материализира във файловата организация на проекта.

\subsection{Обща организация на проекта}

Проектът е структуриран като стандартно Maven приложение. В кореновата директория се намират \texttt{pom.xml}, Spring Boot wrapper файловете, директорията \texttt{src} със сорс код, директорията \texttt{target} с генерирани артефакти, както и директорията \texttt{doc}, в която е разположена настоящата дипломна документация. Това разделение е важно, защото позволява да се разграничат поне три различни аспекта на системата:

\begin{itemize}[leftmargin=*]
    \item \textbf{изходен код и ресурси} -- реализирани в \texttt{src};
    \item \textbf{резултат от компилация и копиране на ресурси} -- реализирани в \texttt{target};
    \item \textbf{академична документация} -- реализирана в \texttt{doc}.
\end{itemize}

Тази подредба е полезна не само от практическа гледна точка, но и за дипломното описание, тъй като ясно отделя инженерния продукт от неговото документиране.

\begin{table}[H]
    \centering
    \caption{Ключови директории в проекта и тяхното предназначение}
    \label{tab:project-directories}
    \begin{tabularx}{\textwidth}{L{3.5cm}Y}
        \toprule
        Директория или файл & Предназначение \\
        \midrule
        \texttt{pom.xml} & Описва зависимостите, версията на Java и build конфигурацията \\
        \texttt{src/main/java} & Съдържа основния Java код на приложението \\
        \texttt{src/main/resources/templates} & Съдържа Thymeleaf шаблоните \\
        \texttt{src/main/resources/static} & Съдържа статични ресурси като CSS и JavaScript \\
        \texttt{src/main/resources/application.properties} & Съдържа локалната конфигурация за база данни, mail и JPA \\
        \texttt{target/classes} & Съдържа компилирани класове и копирани ресурси след build \\
        \texttt{target/test-classes} & Съдържа тестови артефакти при наличие на компилирани тестове \\
        \texttt{doc} & Съдържа LaTeX документацията и библиографията на проекта \\
        \bottomrule
    \end{tabularx}
\end{table}

Както се вижда от таблица \ref{tab:project-directories}, проектът следва добре позната и разбираема структура. Това е съществено предимство, защото разработчик, който е работил със Spring Boot и Maven, може бързо да се ориентира в него.

\subsection{Описание на зависимостите в \texttt{pom.xml}}

Файлът \texttt{pom.xml} е основният декларативен център на Java проекта. В него са зададени родителският Spring Boot артефакт, версията на Java и списъкът със зависимости. Анализът на файла показва, че системата използва сравнително типичен набор от библиотеки за Spring базирано уеб приложение.

\begin{table}[H]
    \centering
    \caption{Основни зависимости и роля в проекта}
    \label{tab:maven-dependencies}
    \begin{tabularx}{\textwidth}{L{4.5cm}YY}
        \toprule
        Зависимост & Функция в приложението & Практическо значение \\
        \midrule
        \texttt{spring-boot-starter-web} & Уеб инфраструктура и MVC маршрути & Осигурява основата на HTTP слоя \cite{springboot,springmvc} \\
        \texttt{spring-boot-starter-thymeleaf} & Сървърно рендирани HTML шаблони & Интегрира визуалния слой с Spring MVC \cite{thymeleaf} \\
        \texttt{spring-boot-starter-data-jpa} & Persistence слой чрез JPA & Улеснява работата с релационната база \cite{springdatajpa} \\
        \texttt{mysql-connector-java} & JDBC драйвер за MySQL & Позволява реална връзка с базата данни \cite{mysql} \\
        \texttt{spring-boot-starter-validation} & Валидация на входни данни & Използва се в binding моделите \cite{hibernatevalidator} \\
        \texttt{modelmapper} & Преобразуване между модели & Намалява ръчния код при mapping \cite{modelmapper} \\
        \texttt{spring-context-support} и \texttt{javax.mail} & Поддръжка за e-mail интеграция & Използват се при потвърждение на регистрацията \cite{springmail} \\
        \texttt{spring-security-crypto} & Хеширане на пароли & Осигурява \texttt{Pbkdf2PasswordEncoder} \cite{springsecuritycrypto} \\
        \texttt{spring-boot-starter-test} & Тестова инфраструктура & Осигурява база за бъдещи тестове \\
        \bottomrule
    \end{tabularx}
\end{table}

Таблица \ref{tab:maven-dependencies} показва, че изборът на зависимости е в общия случай добре съобразен с целите на проекта. Той покрива всички ключови нужди: уеб слой, шаблони, persistence, валидация, mail интеграция и password hashing.

Въпреки това анализът на \texttt{pom.xml} разкрива и няколко важни наблюдения. На първо място, зависимостта \texttt{spring-security-crypto} присъства два пъти с различни версии. На второ място, \texttt{spring-boot-starter-data-jpa} е декларирана с версия \texttt{LATEST}, което не е добра практика в устойчив проект, защото може да доведе до по-трудно възпроизводими build резултати. На трето място, някои security зависимости са коментирани, което подсказва, че проектът първоначално е бил разглеждан и като кандидат за по-богата security интеграция, но впоследствие е останал при по-опростен модел.

\subsection{Версия на платформата и езика}

Проектът използва Spring Boot 2.6.2 и Java 11. Това е разумна комбинация за учебно-приложна разработка. Java 11 е дългосрочно поддържана версия на езика и осигурява стабилна основа, а Spring Boot 2.6.2 предлага достатъчно зрял и добре документиран стек за уеб приложения. Тази комбинация гарантира съвместимост с използваните библиотеки и прави проекта лесен за стартиране в типична локална разработваческа среда.

\subsection{Конфигурация на приложението}

Операционното поведение на приложението се управлява основно от файла \texttt{src/main/resources/application.properties}. В него са зададени параметрите за логване, връзка към MySQL, JPA поведението, както и SMTP конфигурацията за електронната поща.

Конфигурацията на базата данни използва MySQL драйвер и URL към база \texttt{onlineShop}, като е включена опция за автоматично създаване на базата при нужда. Hibernate е настроен с \texttt{ddl-auto=update}, което означава, че схемата може да бъде синхронизирана с entity модела при стартиране. Това решение е удобно за локална разработка, защото намалява необходимостта от отделни миграционни скриптове в ранна фаза на проекта.

От operational гледна точка е важно и използването на \texttt{spring.jpa.open-in-view=false}. Това показва съзнателен избор да не се оставя persistence контекстът отворен през визуалния слой. Макар проектът да не използва сложни lazy loading сценарии в интерфейса, тази настройка е добра практика, защото намалява вероятността от неочаквани достъпи до базата по време на визуализация.

\subsection{Конфигурация на електронната поща}

Файлът \texttt{application.properties} съдържа и настройките за \texttt{spring.mail.host}, порт, SMTP автентикация, TLS и потребителско име. Тези параметри позволяват изпращане на потвърждаващ e-mail към новорегистрирани потребители. От функционална гледна точка това е значимо предимство, защото системата преминава отвъд чисто локалните CRUD операции и демонстрира интеграция с външен инфраструктурен компонент.

Наред с това обаче тук се вижда и едно от най-важните operational ограничения на проекта: чувствителни данни като парола за e-mail акаунт и данни за достъп до базата са записани директно в конфигурационен файл. Това е приемливо за локална учебна среда, но не следва да бъде допускано в production режим. По-зрял подход би включвал използване на environment variables, секретно хранилище или отделни профили за различни среди.

\subsection{Build процес и Maven жизнен цикъл}

Maven \cite{maven} осигурява стандартизиран жизнен цикъл за компилация, тестване и пакетиране. Дори когато проектът не използва сложни custom plugin сценарии, самото наличие на Maven build носи няколко преимущества:

\begin{itemize}[leftmargin=*]
    \item централизирано управление на зависимостите;
    \item възпроизводим build процес на различни машини;
    \item ясна конвенция за разположение на сорс код и ресурси;
    \item лесна интеграция със Spring Boot Maven plugin за стартиране и пакетиране.
\end{itemize}

В рамките на настоящия проект build процесът води до попълване на директорията \texttt{target}, където се намират \texttt{classes}, \texttt{test-classes}, както и стандартни поддиректории за генерирани сорсове. Най-важната от тях е \texttt{target/classes}, в която се копират както компилираните \texttt{.class} файлове, така и ресурсите от \texttt{src/main/resources}, включително HTML шаблоните.

\subsection{Съдържание и роля на \texttt{target}}

От наличната структура се вижда, че \texttt{target} съдържа поне следните поддиректории: \texttt{classes}, \texttt{generated-sources}, \texttt{generated-test-sources} и \texttt{test-classes}. Това е очакван резултат от Maven build процес. Практическото значение на тази директория е двойно.

Първо, тя съдържа реално изпълнимите артефакти на приложението след компилация. Второ, тя предоставя добра възможност за проверка дали ресурсите от \texttt{templates} и \texttt{static} са коректно копирани и достъпни за изпълнимата среда. В конкретния проект например шаблоните от \texttt{src/main/resources/templates} са налични и в \texttt{target/classes/templates}, което потвърждава правилното участие на ресурсите в build процеса.

\subsection{Практическа среда за изпълнение}

За локално стартиране на приложението са необходими поне следните предпоставки:

\begin{itemize}[leftmargin=*]
    \item налична Java 11 среда;
    \item работещ MySQL сървър;
    \item валидни данни за достъп до базата, съобразени с \texttt{application.properties};
    \item достъпна e-mail конфигурация, ако се тества регистрационният поток;
    \item Maven или използване на включения wrapper за стартиране на проекта.
\end{itemize}

Тази operational рамка е характерна за Spring Boot приложения и не изисква специализирана инфраструктура. Именно това е още едно предимство на избрания стек: проектът може да бъде разработван и демонстриран в сравнително стандартна среда без значителни допълнителни инструменти.

\subsection{Операционни рискове и възможности за подобрение}

Анализът на конфигурацията позволява да бъдат формулирани няколко operational риска:

\begin{itemize}[leftmargin=*]
    \item чувствителни параметри са съхранени в изходния код вместо във външна конфигурация;
    \item използването на \texttt{LATEST} в зависимост създава риск за повторяемостта на build процеса;
    \item наличието на дублирани зависимости подсказва необходимост от по-строга dependency хигиена;
    \item липсват ясно разграничени профили за development и production среда;
    \item отсъстват миграции на схемата чрез инструмент като Flyway или Liquibase.
\end{itemize}

Въпреки това operational картината на системата остава напълно достатъчна за целите на дипломната разработка. Проектът е стартиращ, конфигурируем и използва стандартни Java инструменти за компилация и изпълнение. Това го прави едновременно практичен и методологично удобен за анализ.

\subsection{Обобщение}

Структурата на проекта, Maven конфигурацията и operational средата показват, че разглежданата система е реализирана по дисциплиниран начин, следващ конвенциите на Spring Boot и Maven. Наличието на ясно разграничени директории \texttt{src} и \texttt{target}, добре подбран стек от зависимости и централна конфигурация в \texttt{application.properties} осигурява добра основа за локално стартиране и анализ на приложението.

Същевременно operational прегледът разкрива и важни области за бъдещо професионализиране: управление на тайни, dependency консолидация, профили за среди и по-строг контрол върху persistence миграциите. Именно тези наблюдения превръщат настоящата глава в съществена част от цялостната техническа документация, а не само в инвентарен списък на файловете.
